% Jacob Neumann

% DOCUMENT CLASS AND PACKAGE USE
\documentclass[aspectratio=169, handout]{beamer}

% Establish the colorlambda boolean, to control whether the lambda is solid color (true), or the same as the picture (false)
\newif\ifcolorlambda
\colorlambdafalse % DEFAULT: false

% Use auxcolor for syntax highlighting
\newif\ifuseaux
\useauxfalse % DEFAULT: false

% Color settings
\useauxtrue

% LIGHT MODE TOGGLE: flip \lightmodetrue for a projector-friendly light theme
% (dark text on a near-white canvas). The base palette below switches on it, and
% the accent colors further down (ntColor, part*, code*) get darker variants too.
\newif\iflightmode
\lightmodefalse % DEFAULT: dark mode

\iflightmode
  % --- light theme: dark ink on a soft near-white page ---
  \newcommand{\auxColor}{9A6B12}     % darker gold (readable on light)
  \newcommand{\presentColor}{0C4033} % deep green accent
  \newcommand{\bgColor}{F4F1E8}      % soft off-white canvas
  \newcommand{\darkBg}{8b98ad}
  \newcommand{\textColor}{1A2420}    % near-black ink
\else
  % --- dark theme (original) ---
  \newcommand{\auxColor}{CC9A49}     % the color of note boxes and stuff
  \newcommand{\presentColor}{01180E} % the primary color of the slide borders
  \newcommand{\bgColor}{011C12}      % the color of the background of the slide
  \newcommand{\darkBg}{8b98ad}
  \newcommand{\textColor}{E7CEA4}
\fi
\newcommand{\lambdaColor}{\auxColor}

\colorlambdatrue

\usepackage{comment} % comment blocks
\usepackage{soul} % strikethrough
\usepackage{listings} % code
\usepackage{makecell}
\usepackage{transparent}
\usepackage{proof} % PL inference rules: \infer{conclusion}{premises}

\setbeamertemplate{itemize items}[circle]
% \setbeameroption{show notes on second screen=right}

\usepackage{lectureSlides}

% Serif body text. lectureSlides loads Inter (sans) as the default family;
% point the roman family at Crimson and make the document default roman so
% canvas body text is serif.
\usepackage[T1]{fontenc}
\usepackage{crimson}              % sets \rmdefault to the Crimson family
\renewcommand{\familydefault}{\rmdefault}

% The inter package (loaded by lectureSlides) drives math italic letters through
% a `pureletters` symbol font that follows the text roman family. Once crimson
% repointed \rmdefault, math letters rendered in Crimson and clashed with the CM
% symbols. Redeclare `pureletters` back to Computer Modern math italic so math
% stays in CM while body text stays Crimson. Must be \AtBeginDocument: inter
% installs `pureletters` at begin-document, after the preamble.
\AtBeginDocument{\DeclareSymbolFont{pureletters}{OML}{cmm}{m}{it}}

% Override ONLY the frametitle bar colors for this file (must come after
% lectureSlides, which defines the frametitle template + theme). Change the
% hex to recolor the title bar; leave the rest of the theme untouched.
\definecolor{frametitleBar}{HTML}{0C4033}
\setbeamercolor{frametitle}{fg=auxColor, bg=frametitleBar}


% The custom frametitle template (lectureSlides.sty) sets a font but never
% applies the frametitle fg, so beamer's fg= gets dropped. Re-inject it.
\addtobeamertemplate{frametitle}{\usebeamercolor[fg]{frametitle}\color{fg}}{}

% Primary body text color on the canvas. The dark bgColor makes the default
% black text unreadable. Note \textColor is just a macro holding a hex string;
% xcolor/beamer's fg= needs a *registered* color name, so define one first.
\definecolor{textColor}{HTML}{\textColor}
\setbeamercolor{normal text}{fg=textColor}

% TOC entries inherit from `structure` -> presentColor (near-black here), so the
% Table of Contents was invisible. Color the toc entries explicitly.
\setbeamercolor{section in toc}{fg=auxColor}
\setbeamercolor{subsection in toc}{fg=textColor}

% itemize/enumerate markers also inherit from `structure` (near-black), so the
% bullets were invisible on the dark canvas. Color the markers gold.
\setbeamercolor{itemize item}{fg=auxColor}
\setbeamercolor{itemize subitem}{fg=auxColor}
\setbeamercolor{itemize subsubitem}{fg=auxColor}
\setbeamercolor{enumerate item}{fg=auxColor}

% Footnotes inherit from `structure` too, so the footnote text was near-black and
% invisible (only its separator rule showed). \setbeamercolor alone didn't take,
% so force the color directly in the template around \insertfootnotetext.
\setbeamercolor{footnote}{fg=textColor}
\setbeamercolor{footnote mark}{fg=auxColor}
\setbeamertemplate{footnote}{%
  \parindent 1em\noindent%
  \raggedright
  \hbox to 1.8em{\hfil{\color{auxColor}\insertfootnotemark}}%
  {\color{textColor}\insertfootnotetext}\par%
}

% Grammar markup: distinguish nonterminals from terminals by color AND font.
%   \nt{...} nonterminal (a symbol that expands further) -> bright gold, the
%            regular math italic (NOT bold: bold-italic e was muddy/hard to read)
%   \op{...} terminal OPERATOR (+ - * /) -> cream typewriter, kept as a \mathbin
%            so it retains the normal binary-operator spacing
%   \tm{...} terminal literal (numbers, parens, <integer literal>) -> cream
%            typewriter, so program text reads like code (\texttt-style)
% Coloring a math sub-expression with {\color..#1} turns it into an ordinary
% atom and strips operator spacing; \op uses \mathbin to restore it.
% teal nonterminal accent: a bright pastel for dark mode, a deeper teal for light
\iflightmode
  \definecolor{ntColor}{HTML}{0E7C6E}  % deep teal (readable on light)
  \definecolor{linkColor}{HTML}{2A5DB0} % link: deep blue on light
\else
  \definecolor{ntColor}{HTML}{5FD0C0}  % bright teal (for dark)
  \definecolor{linkColor}{HTML}{7FB2E8} % link: soft sky-blue (not garish on dark)
\fi
% \link{url}{text}: a hyperlink styled in the (toggle-aware) link color.
% Plain \url with \color{blue} is too harsh on the dark canvas.
\newrobustcmd{\link}[2]{\href{#1}{\color{linkColor}\underline{#2}}}
% \nt is mode-agnostic (\ensuremath) so it works bare in math AND inside the
% text-mode \gline below. Regular weight (bold was too fat); distinction from
% the cream terminals comes from color + the terminals being typewriter.
\newcommand{\nt}[1]{\ensuremath{{\color{ntColor}#1}}}
% \mathtt does NOT restyle digits/operators (they come from the symbol fonts),
% so terminals must use real text-mode \texttt via \text{} to look like code.
\newcommand{\op}[1]{\mathbin{\text{\color{textColor}\texttt{#1}}}}
\newcommand{\tm}[1]{\text{\color{textColor}\texttt{#1}}}

% \gline{...}: terminal-by-default wrapper. Sets the ambient font to cream
% typewriter (text mode), so bare characters ( ) + * / digits spaces all render
% as terminals with NO markup. Only nonterminals need wrapping in \nt{...}, which
% pops into gold math. \gto is the derivation arrow for use inside \gline.
% \text{} forces text mode so \gline works even inside an align*/array cell.
\newcommand{\gline}[1]{\text{\ttfamily\color{textColor}#1}}
\newcommand{\gto}{\ensuremath{\;\Longrightarrow\;}}
% \kw{...}: a terminal KEYWORD inside \gline -> gold typewriter (matches the
% gold keyword highlighting that \code gives). Everything else in \gline stays
% cream; \nt{...} stays teal.
\newcommand{\kw}[1]{{\color{auxColor}#1}}
% \rname{...}: an inference-rule name for \infer's optional [..] arg -- small
% caps, gold, so rule labels read as a consistent set. Usage:
%   \infer[\rname{Add}]{conclusion}{premises}
\newcommand{\rname}[1]{\text{\small\color{auxColor}\textsc{#1}}}
% \ty{...}: a type name (int, bool, ...) in code/typewriter style, like \tm but
% semantically "a type". Body is in math mode (\ensuremath) so subscripts like
% t_1 work; \mathtt gives the typewriter look for letters. For compound types
% use \op for the connective, e.g. \ty{t_1} \op{*} \ty{t_2}.
\newcommand{\ty}[1]{\ensuremath{{\color{textColor}\mathtt{#1}}}}

% fact_proto part-highlighting: three semantic colors to mark the condition,
% the base case, and the recursive case, then referenced in the explanation.
\iflightmode
  \definecolor{partCond}{HTML}{1763B0}  % blue  (deeper, for light)
  \definecolor{partBase}{HTML}{2E8B40}  % green (deeper, for light)
  \definecolor{partRec}{HTML}{B03C86}   % pink  (deeper, for light)
\else
  \definecolor{partCond}{HTML}{6FB7F0}  % condition  -> blue
  \definecolor{partBase}{HTML}{9FE0A0}  % base case  -> green
  \definecolor{partRec}{HTML}{E59BC8}   % recursive case -> pink
\fi
% (\prec is a built-in relation symbol, so the recursive-case macro is \precase)
\newcommand{\pcond}[1]{{\color{partCond}#1}}
\newcommand{\pbase}[1]{{\color{partBase}#1}}
\newcommand{\cl}[1]{{\color{partCond}#1}}
\newcommand{\precase}[1]{{\color{partRec}#1}}

% \lam{x}: a lambda binder "lambda x." with a TIGHT dot. A bare \lambda x. gives
% the period \mathpunct spacing (a wide gap after the dot); {.} makes the dot an
% ordinary atom and \mkern adds just a hair of space before the body. Usage:
%   \lam{f} \lam{x} f x   ==>   lambda f. lambda x. f x
\newcommand{\lam}[1]{\lambda #1{.}\mkern2mu}

% \Y: the Y combinator, in a poppy magenta so it stands out as the "magic" term.
% Use it (in math mode) anywhere the mathematical Y appears: \Y, \Y \; f, etc.
\definecolor{yColor}{HTML}{F25CA2}
% \newrobustcmd so \color survives inside moving args (e.g. \defBox does an
% \edef, where a fragile \color would break).
\newrobustcmd{\Y}{{\color{yColor}\texttt{Y}}}

% \hlt{NAME}: a high-level lambda-calculus "widget" term (fact, IF, EQ, MUL,
% PRED, ...) in gold typewriter, so the named building blocks stand out from
% plain variables. Robust so it survives \defBox/moving args.
\newrobustcmd{\hlt}[1]{{\color{auxColor}\texttt{#1}}}

% Big Bird subtyping diagram for the "Why semantics?" section: an array of dogs
% where muppetify drops a BigBird into slot 0, then barkify calls .bark() on all
% -- every dog says woof, but the bird (slot 0) can't (???).
\definecolor{birdRed}{HTML}{E0564F}
\usetikzlibrary{positioning,arrows.meta}
\newcommand{\birddiagram}{%
\begin{tikzpicture}[
   font=\ttfamily\footnotesize,
   cell/.style={draw=textColor, minimum width=1.05cm, minimum height=0.95cm},
   dog/.style={cell, text=ntColor},
   bird/.style={cell, draw=birdRed, fill=bgColor, text=birdRed, very thick},
   % gaps live here -- edit these instead of absolute coordinates:
   % woof/??? sit just below the index labels (i##1), which sit below the cells:
   sound/.style={below=3pt of i##1},  % ## because this is inside a \newcommand body
]
  % the array: bird in slot 0, dogs in 1..5
  \node[bird] (c0) at (0,0) {bird};
  \foreach \i in {1,...,5} { \node[dog] (c\i) at (\i*1.15, 0) {dog}; }
  \foreach \i in {0,...,5} { \node[font=\tiny\ttfamily, text=textColor, below=1pt of c\i] (i\i) {[\i]}; }
  % muppetify drops Big Bird into slot 0 (label sits just above the array)
  \node[text=birdRed, font=\scriptsize\bfseries, above=12pt of c0] (mup) {muppetify: animals[0] = BigBird};
  \draw[-{Stealth[length=2mm]}, birdRed, thick] (mup) -- (c0);
  % barkify: .bark() on each -- woof under the dogs, ??? under the bird
  \foreach \i in {1,...,5} { \node[text=auxColor, font=\tiny, sound=\i] (w\i) {woof}; }
  \node[text=birdRed, font=\small\bfseries, sound=0] {???};
  % barkify caption below the woof row, spanning toward the middle
  \node[text=auxColor, font=\scriptsize, below=8pt of w3] {barkify: call \texttt{.bark()} on each};
\end{tikzpicture}}

% Code blocks: the shared 15150code listings style assumes a LIGHT slide (light
% gray panel, dark-ish text). On this dark deck the code text inherited cream and
% vanished on the light box. Re-issue the style with a DARK panel + light text +
% dark-tuned syntax colors. (Overrides lectureSlides; shared .sty untouched.)
\iflightmode
  \definecolor{codeBg}{HTML}{ECEAE0}      % light panel
  \definecolor{codeKeyword}{HTML}{9A6B12} % darker gold
  \definecolor{codeString}{HTML}{8A6A00}  % darker warm gold
  \definecolor{codeComment}{HTML}{5E7468} % muted green-gray (darker)
  \definecolor{codeNumber}{HTML}{6B7570}  % gray
\else
  \definecolor{codeBg}{HTML}{0A2E22}      % panel: a touch lighter than bgColor
  \definecolor{codeKeyword}{HTML}{CC9A49} % keywords -> gold (auxColor); teal is
                                          % reserved for nonterminals
  \definecolor{codeString}{HTML}{E7B95C}  % strings  -> warm gold
  \definecolor{codeComment}{HTML}{7E9C8E} % comments -> muted green-gray
  \definecolor{codeNumber}{HTML}{B8C4BE}  % line numbers -> light gray
\fi
\lstdefinestyle{15150code}{
    basicstyle=\ttfamily\color{textColor},
    numberstyle=\tiny\ttfamily\color{codeNumber},
    backgroundcolor=\color{codeBg},
    stringstyle=\color{codeString},
    keywordstyle=\color{codeKeyword},
    commentstyle=\color{codeComment},
    numbers=left,
    frame=shadowbox,
    rulesepcolor=\color{codeBg},
    linewidth=\textwidth,
    columns=fixed,
    tabsize=2,
    xleftmargin=\MaxSizeOfLineNumbers,
    breaklines=true,
    keepspaces=true,
    showstringspaces=false,
    captionpos=b,
    autogobble=true,
    mathescape=true,
    literate={~}{{$\thicksim$}}1
}

% Tint the top-right lambda. white_lambda.png is a WHITE silhouette whose alpha
% channel IS the glyph mask. Rather than fight PDF blend modes (which tint the
% bounding box), we bake a recolored copy once (RGB = tint, alpha preserved) and
% include that. The tinted asset is generated by tint_lambda.py; rerun it to
% change the color.
\newcommand{\tintedlambda}{%
  \includegraphics[width=0.6cm]{aux_lambda.png}%
}

% Re-point ONLY the image in the frametitle template to the tinted version.
\setbeamertemplate{frametitle}{%
  \nointerlineskip
  \begin{beamercolorbox}[sep=0.2cm,ht=1.3em,wd=\paperwidth]{frametitle}
    \vbox{}\vskip-2ex%
    \usebeamercolor[fg]{frametitle}\color{fg}%
    \intermedium\fontsize{10.5}{12.5}\selectfont
    \strut\insertframetitle\strut
    \hfill
    \raisebox{-2mm}{\tintedlambda}%
    \vskip-1ex%
  \end{beamercolorbox}%
  % Gold rule flush under the title bar, separating it from the canvas.
  % Reuse a full-width beamercolorbox (same mechanism as the bar above) so the
  % rule starts at the exact same left edge — no manual margin math. sep=0
  % removes padding; the box bg is the gold itself.
  \nointerlineskip\vskip-1pt
  \begin{beamercolorbox}[sep=0pt,wd=\paperwidth]{}%
    {\color{auxColor}\rule{\paperwidth}{0.4pt}}%
  \end{beamercolorbox}%
}

% Mirror of the title-bar rule, but ABOVE the bottom (footline) bar. Prepended
% to the footline template so it sits directly atop the bar; same full-width
% beamercolorbox mechanism for identical left-edge alignment.
\addtobeamertemplate{footline}{%
  \begin{beamercolorbox}[sep=0pt,wd=\paperwidth]{}%
    {\color{auxColor}\rule{\paperwidth}{0.4pt}}%
  \end{beamercolorbox}%
  \nointerlineskip
}{}

% Gold side decals on every regular slide. Drawn in the `background` template
% (runs behind content on every frame) with TikZ anchored to the page edges; the
% right copy is the left one mirrored via xscale=-1. Skipped on [plain] frames
% (title page + \sectionSlide dividers) via beamer's \ifbeamer@plainframe.
%
% side_decal.png is a tall vertical strip, so size it by HEIGHT and run it down
% the full page edge. Tunables:
%   \decalheight — rendered height of one decal (default: full paper height)
%   \decalinset  — horizontal gap from the paper edge
\usetikzlibrary{calc}
\newcommand{\decalimg}{side_decal2_gold} % keyed + flooded to auxColor gold, so
                                         % it keeps its tint even when faded
\newlength{\decalheight}\setlength{\decalheight}{\paperheight}
% negative = push decals further toward (partly off) the page edge
\newlength{\decalinset}\setlength{\decalinset}{-2.5mm}
% \decalopacity: 1 = solid, lower = fainter so the ornaments read as faint
% impressions in the background rather than foreground pieces. Applied via the
% TikZ `opacity` node option below (TikZ-native; \transparent breaks inside a
% TikZ node path).
\newcommand{\decalopacity}{0.66}
\newcommand{\sidedecal}{\includegraphics[height=\decalheight]{\decalimg}}
% The template body runs in the document body, where @ is not a letter, so
% \ifbeamer@plainframe can't be used directly. Capture the test into an @-free
% macro now (under \makeatletter): \ifplainframe{<skip>}{<draw>}.
\makeatletter
\newcommand{\ifplainframe}{\ifbeamer@plainframe \expandafter\@firstoftwo \else \expandafter\@secondoftwo \fi}
\makeatother
% The texture/decals/border are baked DARK-green assets, so in light mode we
% suppress all of them and let the light bgColor canvas show through.
\addtobeamertemplate{background}{%
  \iflightmode\else
  \ifplainframe{%
    % PLAIN frames (title, section dividers, break, poll): a full-frame
    % celestial border (16:9, white keyed out) filling the whole slide, plus the
    % mottled texture underneath. Same opacity as the side decals.
    \begin{tikzpicture}[remember picture, overlay]
      \node[anchor=center] at (current page.center)
        {\includegraphics[width=\paperwidth,height=\paperheight]{bg_texture}};
      \node[anchor=center, opacity=\decalopacity] at (current page.center)
        {\includegraphics[width=\paperwidth,height=\paperheight]{celestial_border2_gold}};
    \end{tikzpicture}%
  }{%
    \begin{tikzpicture}[remember picture, overlay]
      % Subtle mottled texture filling the whole page, UNDERNEATH the decals.
      % Generated by make_bg_texture.py; tinted to bgColor.
      \node[anchor=center] at (current page.center)
        {\includegraphics[width=\paperwidth,height=\paperheight]{bg_texture}};
      % left decal, against the left edge (opacity fades it into the background)
      \node[anchor=west, opacity=\decalopacity] at ([xshift=\decalinset]current page.west)  {\sidedecal};
      % right decal: mirror the IMAGE with \reflectbox (flipping the node's
      % transform instead would move it off-page), anchored at the right edge.
      \node[anchor=east, opacity=\decalopacity] at ([xshift=-\decalinset]current page.east) {\reflectbox{\sidedecal}};
    \end{tikzpicture}%
  }%
  \fi
}{}

% Keep body text clear of the decals by widening the L/R text margins to the
% decal width plus a little breathing room.
\setbeamersize{text margin left=1.2cm, text margin right=1.6cm}

% The shared title-page template positions the full-bleed title image relative
% to a \textwidth-centered \makebox, so changing the text margins above shifted
% it sideways. Re-anchor the image to the absolute page center so it fills the
% slide regardless of margins.
\setbeamertemplate{title page}{%
  \begin{tikzpicture}[remember picture, overlay]
    \node at (current page.center) {\includegraphics[width=\paperwidth,height=\paperheight]{title.png}};
  \end{tikzpicture}%
}

% Plain-slide text fixes. The shared \sectionSlide / \lecturePassword color
% their text with presentColor, which is near-black on this dark deck and so
% invisible. Override them locally (shared .sty untouched) to use gold.
\renewrobustcmd{\sectionSlide}[2]{%
  \section{#2}%
  \begin{frame}[plain]
    \begin{center}
      \huge \color{auxColor} \intermedium{#1 - #2}
    \end{center}
  \end{frame}%
}
\renewrobustcmd{\lecturePassword}[1]{%
  \begin{frame}[plain]
    \begin{center}
      {\huge \color{auxColor} \intermedium{Post-lecture attendance poll}}

      \vspace{10pt}

      \color{auxColor} \intermedium{Password: #1}
    \end{center}
  \end{frame}%
}

%%%%%%%%%%%%%%%%%%%%%%%%%%%%%%%%%%%%%%%%%| <----- Don't make the title any longer than this
\title{Lambda Calculus} % TODO
\subtitle{The foundation of functional programming} % TODO
\date{22 June 2026} % TODO
\author{Brandon Wu} % TODO

\graphicspath{ {./img/} }
% DONT FORGET TO PUT [fragile] on frames with codeblocks, specs, etc.
    %\begin{frame}[fragile]
    %\begin{codeblock}
    %fun fact 0 = 1
    %  | fact n = n * fact(n-1)
    %\end{codeblock}
    %\end{frame}

% INCLUDING codefile:
    % 1. In some file under code/NN (where NN is the lecture id num), include:
%       (* FRAGMENT KK *)
%           <CONTENT>
%       (* END KK *)

%    Remember to not put anything on the same line as the FRAGMENT or END comment, as that won't be included. KK here is some (not-zero-padded) integer. Note that you MUST have fragments 0,1,...,KK-1 defined in this manner in order for fragment KK to be properly extracted.
    %  2. On the slide where you want code fragment K
            % \smlFrag[color]{KK}
    %     where 'color' is some color string (defaults to 'white'. Don't use presentColor.
%  3. If you want to offset the line numbers (e.g. have them start at line 5 instead of 1), use
            % \smlFragOffset[color]{KK}{5}

\begin{document}

% Make it so ./mkWeb works correctly
\ifweb
\renewcommand{\pause}{}
\fi

\setbeamertemplate{itemize items}[circle]

% SOLID COLOR TITLE (see SETTINGS.sty)
{
\begin{frame}[plain]
\colorlambdatrue
\titlepage
\end{frame}
}

\begin{frame}[fragile]
\frametitle{Lesson Plan}

\tableofcontents
\end{frame}


\sectionSlide{1}{The Origin of Programming Languages}

\begin{frame}[fragile]
\frametitle{What is Programming Language Theory?}

What is programming language theory?

\pause
\vspace{\fill}

Programming language theory is the study of the design, implementation,
and analysis of programming languages.

\pause
\vspace{\fill}

We are interested in answering questions like:
\begin{itemize}
  \item Can we design a programming language that makes common
  idioms easy to express?
  \item Is it possible for a program in a given language to
  run into undefined behavior?
  \item Can we construct a program that violates the rules of
  type safety?
  \item How do we design a programming language such that we can
  compile it efficiently?
\end{itemize}
\end{frame}

\begin{frame}[fragile]
  \frametitle{Syntax and Semantics}

  The way to look at any programming language is as a combination of
  its syntax and semantics.

  \pause
  \vspace{\fill}

  \defBox{}{\, The \term{syntax} of a programming language is how you write
  programs in the language.
  }

  \pause
  \vspace{\fill}

  \defBox{}{\, The \term{semantics} of a programming language is the set of rules
  governing the behavior of programs in the language. }

  \pause
  \vspace{\fill}

  For instance, we might say that the syntax of an SML program could
  be \code{val x = 2}, and it has the semantics of shadowing any
  existing binding of \code{x}.
\end{frame}

\begin{frame}[fragile]
  \frametitle{Describing Syntax: Grammars}

  Formally, we can describe the syntax of a programming language using
  a \term{grammar}, a set of recursive rules that describe
  the form of a valid program in the language.

  \pause
  \vspace{\fill}

  For instance, for a very simple language consisting of only
  arithmetic expressions, we might write the \term{Backus-Naur form (BNF)}
  grammar as follows:

  \[
  \begin{array}{r r l r}
    \nt{e} & ::= & \nt{intlit} & \text{(integer literal)} \\
      &   | & \nt{e} \op{+} \nt{e} & \text{(addition)} \\
      &   | & \nt{e} \op{*} \nt{e} & \text{(multiplication)}\\
      &   | & \nt{e} \op{-} \nt{e} & \text{(subtraction)}\\
      &   | & \nt{e} \op{/} \nt{e} & \text{(division)}\\
  \end{array}
  \]
\end{frame}

\begin{frame}[fragile]
  \frametitle{Grammars: Terminals and Nonterminals}

  Here, we say that \nt{e} is a \term{nonterminal symbol} ranging over the
  \term{sort} of expressions. In layman's terms, $e$ is a placeholder for any
  possible expression, defined by the rules given below.

  \[
  \begin{array}{r r l r}
    \nt{e} & ::= & \nt{intlit} & \text{(integer literal)} \\
      &   | & \nt{e} \op{+} \nt{e} & \text{(addition)} \\
      &   | & \nt{e} \op{*} \nt{e} & \text{(multiplication)}\\
      &   | & \nt{e} \op{-} \nt{e} & \text{(subtraction)}\\
      &   | & \nt{e} \op{/} \nt{e} & \text{(division)}\\
  \end{array}
  \]

  \pause
  \vspace{\fill}

  We call the integer literal rule a \term{terminal rule}, because it
  does not result in any further recursive calls to rules. In other
  words, it cannot be expanded any further.

  \pause
  \vspace{\fill}

  By contrast, the other four rules are \term{non-terminal rules}, because
  they require further expansion using the definition of \nt{e}, an expression.
\end{frame}




\begin{frame}[fragile]
  \frametitle{Validating Syntax}

  Then, equipped with our grammar, we can validate whether a given
  program is syntactically valid by whether or not it can be produced
  by using a finite number of our production rules.

  \pause
  \vspace{\fill}

  Let's suppose that we wanted to determine if the expression
  $\tm{(1 / 2) + (3 * 4)}$ is syntactically valid.

  \pause
  \vspace{\fill}

  To do this, we start from the \textit{start symbol} of our grammar,
  which is $e$\footnote{Here, we are underlining the place where we expanded the rule definition.}.

  \pause
  \vspace{-7.5pt}

  \begin{minipage}[t]{0.4\textwidth}
    \[
    \begin{array}{r r l r}
      \nt{e} & ::= & \nt{intlit} & \text{(IL)} \\
        &   | & \nt{e} \op{+} \nt{e} & \text{(ADD)} \\
        &   | & \nt{e} \op{*} \nt{e} & \text{(MUL)}\\
        &   | & \nt{e} \op{-} \nt{e} & \text{(SUB)}\\
        &   | & \nt{e} \op{/} \nt{e} & \text{(DIV)}\\
    \end{array}
    \]
  \end{minipage}
  \begin{minipage}[t]{0.55\textwidth}
    \begin{align*}
      \nt{e} & \stepsTo \gline{\underline{\nt{e_1} + \nt{e_2}}} & \tag{ADD} \\
               & \stepsTo \gline{\underline{(\nt{e_3} / \nt{e_4})} + \nt{e_2}} & \tag{DIV} \\
               & \stepsTo \gline{(\nt{e_3} / \nt{e_4}) + \underline{(\nt{e_5} * \nt{e_6})}} & \tag{MUL} \\
               & \stepsTo \gline{(\underline{1} / \underline{2}) + (\underline{3} * \underline{4})} & \tag{IL}
    \end{align*}
  \end{minipage}

\end{frame}

\begin{frame}[fragile]
  \frametitle{Grammar of an SML Subset}

  For instance, we might describe a limited subset of the grammar of an SML program as follows:

  \pause

  \[
  \begin{array}{r r l r}
    \nt{e} & ::= & \gline{\kw{fn} \nt{p} \kw{=>} \nt{e}} & \text{(lambda expression)}\\
      &   | & \gline{\nt{e_1} \nt{e_2}} & \text{(application)} \\
      &   | & \gline{x} & \text{(variable)} \\
      &   | & \nt{intlit} & \text{(integer literal)} \\
  \end{array}
  \]

  \pause
  \vspace{\fill}

  \[
  \begin{array}{r r l r}
    \nt{p} & ::= & \nt{id} & \text{(identifier)}\\
      &   | & \gline{(\nt{p_1}, $\dots$, \nt{p_n})} & \text{(tuple, $n \geq 2$)} \\
      &   | & \gline{\nt{id} \nt{p}} & \text{(constructor)} \\
      &   | & \gline{\kw{\_}} & \text{(wildcard)} \\
  \end{array}
  \]

  \pause
  \vspace{\fill}

  \[
  \begin{array}{r r l r}
    \nt{dec} & ::= & \gline{\kw{val} \nt{p} \kw{=} \nt{e}} & \text{(value declaration)}\\
        &   | & \gline{\kw{fun} \nt{id} \nt{args} \kw{=} \nt{e}} & \text{(function declaration)} \\
        &   | & \gline{\nt{dec_1} \nt{dec_2}} & \text{(sequence)} \\
  \end{array}
  \]

\end{frame}

\begin{frame}[fragile]
  \frametitle{Deriving SML Syntax}

  Let's look at how we might derive the syntax of the SML program\footnote{A reminder that SML programs are whitespace
  agnostic, so we omit the newline in the derived program.}
  \begin{codeblock}
    val x = 1
    val _ = fn f => f x
  \end{codeblock}

  \pause
  \vspace{-5pt}

  % Tighten the derivation so the slide leaves room for the footnote above the
  % footline: smaller font + reduced inter-row gap (\jot).
  \small
  \setlength{\jot}{3pt}

  % \begin{minipage}{0.4\textwidth}

  % \[
  % \begin{array}{r r l r}
  %   \nt{e} & ::= & \code{fn } \;\; \nt{p}\;\; \code{ => } \;\; \nt{e} & \text{(LAM)}\\
  %     &   | & \nt{e_1} \;\; \nt{e_2} & \text{(APP)} \\
  %     &   | & \code{x} & \text{(VAR)}
  % \end{array}
  % \]

  % \[
  % \begin{array}{r r l r}
  %   \nt{p} & ::= & \nt{id} & \text{(ID)}\\
  %     &   | & (\nt{p_1}, ..., \nt{p_n}) & \text{(TUP, $n \geq 2$)} \\
  %     &   | & \nt{id} \;\; \nt{p} & \text{(CONS)} \\
  %     &   | & \code{_} & \text{(WILDCARD)} \\
  % \end{array}
  % \]

  % \[
  % \begin{array}{r r l r}
  %   \nt{dec} & ::= & \code{val } \;\; \nt{p} \;\; \code{ = } \;\; \nt{e} & \text{(VAL)}\\
  %       &   | & \code{fun } \;\; \nt{id} \;\; \nt{args} \;\; \code{ = } \;\; \nt{e} & \text{(FUN)} \\
  %       &   | & \nt{dec_1} \; \nt{dec_2} & \text{(SEQ)} \\
  % \end{array}
  % \]
  % \end{minipage}
  % \begin{minipage}{0.55\textwidth}
  \begin{align*}
    \nt{dec} & \stepsTo \gline{\underline{\nt{dec_1} \nt{dec_2}}} & \tag{sequence} \\
              & \stepsTo \gline{\underline{val \nt{p} = \nt{e}} \nt{dec_2}} & \tag{value declaration} \\
              & \stepsTo \gline{val \underline{x} = \underline{1} \nt{dec_2}} & \tag{identifier, integer literal} \\
              & \stepsTo \gline{val x = 1 \underline{val \nt{p} = \nt{e}}} & \tag{value declaration} \\
              & \stepsTo \gline{val x = 1 val \underline{\_} = \underline{fn \nt{p} => \nt{e}}} & \tag{wildcard, lambda expression} \\
              & \stepsTo \gline{val x = 1 val \_ = fn \underline{f} => \nt{e}} & \tag{identifier} \\
              & \stepsTo \gline{val x = 1 val \_ = fn f => \underline{\nt{e_1} \nt{e_2}}} & \tag{application} \\
              & \stepsTo \gline{val x = 1 val \_ = fn f => \underline{f} \underline{x}} & \tag{variable (x2)} \\
  \end{align*}

  \vspace{-10pt}
  % \end{minipage}
\end{frame}

\begin{frame}[fragile]
  \frametitle{Parsing in Practice}

  Usually, we do not actually implement parsers in the way that we have
  described here, where we infinitely expand the grammar until we match
  the target.

  \pause
  \vspace{\fill}

  A detailed treatment of how parsing is implemented is beyond the scope of
  this lecture, and is also generally regarded by programming language
  theorists as an uninteresting, solved part of the problem.

  \pause
  \vspace{\fill}

  We will now turn our attention to the rich problem of the semantics
  of a programming language.
\end{frame}

\sectionSlide{2}{Semantics and Judgements}

\begin{frame}[fragile]
  \frametitle{Why Semantics?}

  Why should we care about programming language theory?

  \vspace{\fill}

  Let's take a foray into the world of Java. Java has a notion of
  \term{subtyping} given by its system of inheritance.

  \begin{codeblock}[language=java]
    class Dog extends Animal {
      ...
    }
  \end{codeblock}

  \vspace{\fill}

  In its simplest form, you can define a class \code{Dog} that
  extends \code{Animal}, which
  means that a \code{Dog} should be admissible anywhere an
  \code{Animal} is allowed.\footnote{Put another way, "ain't no rule
  that says you can't put a Dog where an Animal is expected".}

  \pause
  \vspace{\fill}

  We write this as \code{Dog <: Animal}, meaning that \code{Dog} is a
  subtype of \code{Animal}.
\end{frame}

\begin{frame}[fragile]
  \frametitle{Java's Billion-Dollar Mistake}

  Programming language theorists like to say that Java's type system is
  fundamentally broken\footnote{I have heard these exact words.}.

  \pause
  \vspace{\fill}

  The reason comes from Java's subtyping rules on arrays, which are allowed
  to be generic in their types, like lists in SML. We write that \code{T[]}
  is the type of arrays of elements of type \code{T}.

  \pause
  \vspace{\fill}

  Unfortunately, Java arrays are also \term{covariant}, meaning that if
  \code{B <: A}, then \code{B[] <: A[]}.
\end{frame}

\begin{frame}[fragile]
  \frametitle{A Subtyping Bug}

  However, if this is true, consider the following code:

  { \small
  \begin{codeblock}[language=java]
    public void muppetify(Animal[] animals) {
      animals[0] = new BigBird();
    }

    public void barkify(Dog[] dogs) {
      for (Dog dog : dogs) {
        dog.bark(); // woof
      }
    }

    public void main(String[] args) {
      Dog[] dogs = new Dog[5];
      muppetify(dogs); // dogs can go where animals can!
      barkify(dogs); // but this makes big bird bark :(
    }
  \end{codeblock}
  }

\end{frame}

\begin{frame}[fragile]
  \frametitle{What Went Wrong}

  We see that there \textit{should} be a rule that prevents
  dogs from being used where we expect animals!

  \pause
  \vspace{\fill}

  We see that \code{muppetify} essentially just puts a bird into
  the array of animals, which in \code{main} is \textit{really} an
  array of dogs.

  \pause
  \vspace{\fill}

  This means that we have poor Big Bird in the array of dogs, which
  causes a runtime exception when we try to make him bark.

  \pause
  \vspace{\fill}

  \begin{center}
    \birddiagram
  \end{center}
\end{frame}

\begin{frame}[fragile]
  \frametitle{Why This Matters}

  This might seem like a contrived example or silly trivia, but it is
  really not.

  \pause
  \vspace{\fill}

  Because of this behavior, Java needs to be able to detect at runtime
  whether or not you are doing an unsafe operation, like storing Big Bird
  in an array of dogs\footnote{You only make that mistake once, though.}.

  \pause
  \vspace{\fill}

  This means that arrays actually store additional runtime type information
  to be able to detect this unsafe operation on array load, meaning that
  this actually has a real performance cost.

  \pause
  \vspace{\fill}

  The study of the semantics of programming languages is so that we can
  discover new patterns of writing expressive programs, as well as to
  ensure that our programming languages are \textbf{safe and correct}.
\end{frame}

\begin{frame}[fragile]
  \frametitle{Properties We Care About}

  For instance, we are interested in certain properties of programs like
  \term{progress} and \term{preservation}.

  \pause
  \vspace{\fill}

  \defBox{}{\term{Progress} is the property such that every well-typed expression
  \code{e} is either a value, or there exists \code{e'} such that \code{e} $\stepsTo$ \code{e'}.}

  \pause
  \vspace{\fill}

  \defBox{}{\term{Preservation} is the property that stepping a well-typed expression
  preserves its type. Formally, if \code{e : t} and
  \code{e} $\stepsTo$ \code{e'}, then \code{e' : t}.}

  \pause
  \vspace{\fill}

  How do we prove this? First we need to have a specification for
  what the properties of the programming language are. These are its
  \term{semantics}.
\end{frame}

\begin{frame}[fragile]
  \frametitle{Static and Dynamic Semantics}

  Semantics come in two forms, \term{static} and \term{dynamic}.

  \pause
  \vspace{\fill}

  \defBox{}{\term{Static semantics} are the rules of a programming
  language that govern its properties before being run (e.g. types)
  }

  \pause
  \vspace{\fill}

  \defBox{}{\term{Dynamic semantics} are the rules of a programming
  language that govern its properties during evaluation (how it runs)
  }

  \pause
  \vspace{\fill}

  In programming language theory, the language we use to communicate
  semantics is \term{judgements}.
\end{frame}

\begin{frame}[fragile]
  \frametitle{Judgements}

  \defBox{}{A \term{judgement} is a rule that lets us derive a property
  of a program, given certain assumptions.}

  \pause
  \vspace{\fill}

  In generality, it is notated as:

  % One centered line with both forms and a text "or" between them. (align* with
  % &...&...& would spread the pieces to the margins; "(or)"/conclusions in plain
  % math render as italic letter-products, so they need \text{}.)
  {\large
  \[
    \infer{\mathcal{J}}{\mathcal{J}_1,\; \mathcal{J}_2,\; \dots,\; \mathcal{J}_n}
    \qquad \text{(or)} \qquad
    \infer{\text{conclusion}}{\text{assumptions}}
  \]
  }

  where each $\mathcal{J}$ is itself a judgement.

  \pause
  \vspace{\fill}

  This essentially means that we can derive the fact $\mathcal{J}$,
  if we are given the facts $\mathcal{J}_1, \mathcal{J}_2, \dots, \mathcal{J}_n$.

  \pause
  \vspace{\fill}

  Notably, each $\mathcal{J}_i$ is itself a judgement, which may have assumptions
  that must be fulfilled.

\end{frame}

\begin{frame}[fragile]
  \frametitle{Example Judgements: Days of the Week}

  % A row of inference rules: an array (NOT align*) centers them and keeps even
  % gaps; align* would mis-space the &-separated rules. 2x2 so they fit the width.

  Here are some example judgements:

  \vspace{1em}

  % Two side-by-side blocks, vertically centered against each other ([c], not
  % [t] -- [t] floated the right-hand tree up against the title bar).
  % widths sum to < 1.0 and the trailing % kills the inter-block space; otherwise
  % 0.6 + 0.4 = full width leaves no room and the second minipage wraps below.
  % [c] centers the short left block against the taller right tree, so they read
  % as two tidy side-by-side columns (with [t] the left 2x2 staggered against the
  % 5-deep tree). \hfill + widths < 1.0 keep them on one line.
  \begin{minipage}[c]{0.56\textwidth}
  \[
  \begin{array}{cc}
    \infer{\text{Day $n$ is Sunday}}{\text{Day $n-1$ is Saturday}} &
    \infer{\text{Day $n$ is Monday}}{\text{Day $n-1$ is Sunday}} \\[1.5em]
    \infer{\text{Day $n$ is Tuesday}}{\text{Day $n-1$ is Monday}} &
    \infer{\text{Day $n$ is Wednesday}}{\text{Day $n-1$ is Tuesday}}
  \end{array}
  \]
  \end{minipage}\hfill
  \begin{minipage}[c]{0.4\textwidth}
    \centering
    \scalebox{0.9}{$\infer{\text{Day $n$ is Wednesday}}{
      \infer{\text{Day $n-1$ is Tuesday}}{
        \infer{\text{Day $n-2$ is Monday}}{
          \infer{\text{Day $n-3$ is Sunday}}{
            \text{Day $n-4$ is Saturday}
          }
        }
      }
    }$}
  \end{minipage}

  \pause
  \vspace{\fill}

  On the left, we see the judgements that we are allowed to use, in order
  to derive facts.

  \pause
  \vspace{\fill}

  In plainspeak, if we know that four days ago was Saturday, then we can
  successfully conclude that today is Wednesday, my dudes\footnote{This
  joke might be funnier if it wasn't Monday.}.
\end{frame}

\begin{frame}[fragile]
  \frametitle{Judgements for Programming Languages}

  That's an example of an intuitive judgement, but what about judgements
  that we can use for programming languages?

  \pause
  \vspace{\fill}

  Actually, you've already seen these. Here are some judgements that you
  might see for SML:

  % A single rule -> \[..\], not align*. The trailing & (with nothing after it)
  % was breaking align*. Premises are separated by &; widen the gap between them
  % with \quad / \qquad.
  % \op{...} renders the operator in code/typewriter style AND keeps its binary-
  % operator spacing (it's a \mathbin). \code here would inject an \hbox that
  % breaks the surrounding math spacing.
  \[
    \infer[\rname{S-Add}]{e_1 \op{+} e_2 : \ty{int}}{e_1 : \ty{int} \quad e_2 : \ty{int}}
    \qquad
    \infer[\rname{S-Mul}]{e_1 \op{*} e_2 : \ty{int}}{e_1 : \ty{int} \quad e_2 : \ty{int}}
    \qquad
    \infer[\rname{S-Lit}]{n : \ty{int}}{}
  \]
  \[
    \infer[\rname{S-Tup}]{(e_1, e_2) : \ty{t_1} \op{*} \ty{t_2}}{e_1 : \ty{t_1} \quad e_2 : \ty{t_2}}
  \]

  \pause
  \vspace{\fill}

  We note that the \rname{S-Lit} judgement has no premises, because our
  variable $n$ is quantified over all integer literals. Those require
  no justification to be at type \code{int}.

  \pause
  \vspace{\fill}

  We also typically annotate each judgement with a name so that we can
  write it in the derivation tree so it's clear which judgement we are invoking.

\end{frame}

\begin{frame}[fragile]
  \frametitle{A Typing Derivation}

  Using these judgements, we can derive the fact that the expression
  \code{(1 + 2) * (3 + 4)} is well-typed at type \code{int}:

  \pause
  \vspace{\fill}

  \[
    \infer[\rname{S-Mul}]{(1 \op{+} 2) \op{*} (3 \op{+} 4) : \ty{int}}{
      \infer[\rname{S-Add}]{1 \op{+} 2 : \ty{int}}{
        \infer[\rname{S-Lit}]{1 : \ty{int}}{} \qquad
        \infer[\rname{S-Lit}]{2 : \ty{int}}{}
      }
      \qquad
      \infer[\rname{S-Add}]{3 \op{+} 4 : \ty{int}}{
        \infer[\rname{S-Lit}]{3 : \ty{int}}{} \qquad
        \infer[\rname{S-Lit}]{4 : \ty{int}}{}
      }
    }
  \]

  \pause
  \vspace{\fill}

  If you squint, this is actually no different than the justification
  we've been using all semester as to why certain expressions are well-typed!

  \pause
  \vspace{\fill}

  The great thing about judgements is that they make it very easy for a
  computer to implement what we intuitively understand. It would be very
  easy to write a program to follow these rules and verify if a given
  expression is of type \code{int}, for instance.
\end{frame}

\begin{frame}[fragile]
  \frametitle{Static vs. Dynamic Judgements}

  These judgements we just saw are judgements for static semantics--they
  only govern how we type-check a given program.

  \pause
  \vspace{\fill}

  We also have judgements for dynamic semantics, which tell us how to
  evaluate a given program. We write $e \Downarrow v$ to mean ``$e$ evaluates
  to the value $v$'':
  \[
    \infer[\rname{D-Lit}]{n \Downarrow n}{}
    \qquad
    \infer[\rname{D-Add}]{e_1 \op{+} e_2 \Downarrow n}
                       {e_1 \Downarrow n_1 \quad e_2 \Downarrow n_2 \quad n = n_1 + n_2}
    \qquad
    \infer[\rname{D-Tuple}]{(e_1, e_2) \Downarrow (v_1, v_2)}
                       {e_1 \Downarrow v_1 \quad e_2 \Downarrow v_2}
  \]

  \pause
  \vspace{\fill}

  These three judgements essentially tell us that numbers step to
  themselves, addition takes place after evaluating both operands, and
  a tuple evaluates to the pair of the evaluations of its components.
\end{frame}

\begin{frame}[fragile]
  \frametitle{Dynamic Semantics: Evaluation}

  We can also use these judgements to derive the fact that the expression
  \code{(1 + 2, 3 + 4)} evaluates to the value \code{(3, 7)}:

  \[
    \infer[\rname{D-Tuple}]{(1 + 2, 3 + 4) \Downarrow (3, 7)}{
      \infer[\rname{D-Add}]{1 + 2 \Downarrow 3}{
        \infer[\rname{D-Lit}]{1 \Downarrow 1}{} \quad
        \infer[\rname{D-Lit}]{2 \Downarrow 2}{}
      } \quad
      \infer[\rname{D-Add}]{3 + 4 \Downarrow 7}{
        \infer[\rname{D-Lit}]{3 \Downarrow 3}{} \quad
        \infer[\rname{D-Lit}]{4 \Downarrow 4}{}
      }
    }
  \]

  \pause
  \vspace{\fill }

  In reality, there would be far more rules for evaluating programs.
  Proving a property like progress or preservation would then be a matter
  of using induction to prove that no matter how many of these rules we
  chain together, we will always end up with a program with the desired
  property.
\end{frame}

\begin{frame}[fragile]
  \frametitle{Proving a Safety Property}

  We can sketch out a proof of a safety property for our language:
  if $e : \ty{t}$ and $e \Downarrow v$, then $v : \ty{t}$.

  \pause
  \vspace{\fill}

  This involves a technique called \term{rule induction}. Essentially,
  we will show that through all ways to obtain the judgement $e \Downarrow v$,
  the property holds.

  \pause
  \vspace{\fill}

  % Inline judgements use math mode ($...$) with the \op/\ty macros, NOT \code:
  % \code is verbatim (\lstinline), so \op, $, etc. inside it error out.
  We proceed by rule induction on the judgement $e \Downarrow v$.
  \begin{itemize}
    \item \textbf{Case 1: \rname{D-Lit}} \\
    Here, $e = n$ and $v = n$. By the statics rule \rname{S-Lit}, $n : \ty{int}$,
    so $\ty{t} = \ty{int}$ and indeed $v : \ty{t}$. \\
    \item \textbf{Case 2: \rname{D-Add}} \\
    Here, $e = e_1 \op{+} e_2$ and $v = n$. The only statics rule for
    $e_1 \op{+} e_2$ is \rname{S-Add}, so $\ty{t} = \ty{int}$; and $n : \ty{int}$
    by \rname{S-Lit}. So $v : \ty{t}$.
  \end{itemize}
\end{frame}

\begin{frame}[fragile]
  \frametitle{Safety Proof: The Tuple Case}

  \begin{itemize}
    \item \textbf{Case 3: \rname{D-Tuple}} \\
    Here, $e = (e_1, e_2) \Downarrow (v_1, v_2) = v$, with assumptions
    $e_1 \Downarrow v_1$ and $e_2 \Downarrow v_2$. \\[2pt]
    The only statics rule for $(e_1, e_2)$ is \rname{S-Tup}, so
    $\ty{t} = \ty{t_1} \op{*} \ty{t_2}$ with $e_1 : \ty{t_1}$ and $e_2 : \ty{t_2}$. \\[2pt]
    \pause
    Now we use the IH: since $e_1 : \ty{t_1}$ and
    $e_1 \Downarrow v_1$, the IH gives $v_1 : \ty{t_1}$ and $v_2 : \ty{t_2}$. \\[2pt]
    Then by \rname{S-Tup}, $(v_1, v_2) : \ty{t_1} \op{*} \ty{t_2} = \ty{t}$, as desired.
  \end{itemize}

  \pause
  \vspace{\fill}

  Proofs of this kind can be rather tedious, but can be automated in
  some cases. Usually, when we add more expressive constructs to our
  languages, these proofs become more interesting and complex.
\end{frame}

\sectionSlide{3}{Lambda Calculus and the Y Combinator}

\begin{frame}[fragile]
  \frametitle{Toward the Lambda Calculus}

  So far, we have discussed a bit of the machinery that we use
  to reason about programming languages.

  \pause
  \vspace{\fill}

  These tools are only useful to us insofar as we have access to
  well-defined statics and dynamics that we can write proofs on.

  \pause
  \vspace{\fill}

  For simple expression languages, this is easy, but more realistic
  languages do not always come with simple statics and dynamics. How
  would we even mathematically render the semantics of a large language
  like Java?

  \pause
  \vspace{\fill}

  The answer is the lambda calculus.
\end{frame}

\begin{frame}[fragile]
  \frametitle{All Languages Look Alike}

  If you squint your eyes, all programming languages look pretty
  much the same.

  \pause
  \vspace{\fill}

  \begin{minipage}{0.51\textwidth}
  They all have some combination of:
  \begin{itemize}
    \item variable declarations
    \item function application
    \item conditionals
    \item arithmetic expressions
  \end{itemize}

  \pause
  \vspace{\fill}

  The \term{lambda calculus} is a formal system that is essentially
  the simplest possible programming language. The key observation
  is in observing that we can model any programming language as
  the lambda calculus with some additional special constructs.
  \end{minipage}
  \hfill
  \begin{minipage}{0.47\textwidth}
    \begin{center}
    \includegraphics[width=0.8\textwidth]{pl_same.png}
    \end{center}
  \end{minipage}
\end{frame}

\begin{frame}[fragile]
  \frametitle{The Lambda Calculus: Grammar}

  The lambda calculus has a very simple grammar:

  \[
  \large
  \begin{array}{r r l r}
    \nt{e} & ::= & \lam{\nt{x}} \nt{e} & \text{(lambda abstraction)}\\
      & |   & \nt{e_1} \; \nt{e_2} & \text{(application)} \\
      & |   & \nt{x} & \text{(variable)}
  \end{array}
  \]

  \pause
  \vspace{\fill}

  Here, you should read $\lambda x.e$ as \code{fn x => e},
  an anonymous function that takes in a variable named $x$ and evaluates
  to the expression $e$.

  \pause
  \vspace{\fill}

  That's literally it.
\end{frame}

\begin{frame}[fragile]
  \frametitle{The Lambda Calculus: Dynamics}

  \begin{minipage}{0.33\textwidth}

    The dynamics of the lambda calculus are that there are only functions,
    variables, and function applications.

    \vspace{5pt}

    There are no numbers, no tuples, and no conditionals.
  \end{minipage}
  \hfill
  \begin{minipage}{0.65\textwidth}
    \begin{center}
      \includegraphics[width=0.95\textwidth]{all_lambda_calculus.png}
    \end{center}
  \end{minipage}

  \pause
  \vspace{\fill}

  This might seem strange, because we are used to programming languages
  which have more in-built features. Despite this, the lambda calculus
  remains \term{Turing-complete}, meaning that it can be used to compute
  anything a general purpose computer can.
\end{frame}

\begin{frame}[fragile]
  \frametitle{Beta-Reduction, Formally}

  Formally, the semantics of the lambda calculus are given by the
  relation of \term{$\beta$-equivalence}.
  \vspace{-15pt}
  \[
    \infer[\rname{REFL}]{M \equiv_\beta M}{}
    \qquad
    \infer[\rname{COMM}]{M \equiv_\beta M'}{M' \equiv_\beta M}
    \qquad
    \infer[\rname{TRANS}]{M \equiv_\beta M''}{M \equiv_\beta M' \quad M' \equiv_\beta M''}
  \]
  \[
    \infer[\rname{APP}]{M_1 \; M_2 \equiv_\beta M_1' \; M_2'}{M_1 \equiv_\beta M_1' \quad M_2 \equiv_\beta M_2'}
    \qquad
    \infer[\rname{LAM}]{\lam{x}M \equiv_\beta \lam{x}M'}{M \equiv_\beta M'}
  \]
  \[
    \infer[\rname{BETA}]{(\lam{x}N)\, M \equiv_\beta N[x \mapsto M]\footnote{This means "the term $N$ where we substitute
    $M$ for $x$". There is a bit of a subtlety here having to do with substitution and
    variable capture---refer to {\color{linkColor}\href{https://www.cs.cmu.edu/~rwh/pfpl/supplements/ulc.pdf}{Bob Harper's notes for more}}}}{}
  \]

  Only the last rule is interesting--the other ones just tell us that
  $\beta$-equivalence is an equivalence relation, and recursively it
  can be applied throughout a term.

  \pause
  \vspace{\fill}

  \keyBox{}{\, The last rule is interesting because it essentially tells us the
  semantics of function application.}
\end{frame}


\begin{frame}[fragile]
  \frametitle{Encoding Data with Lambdas}

  This seems ludicrous, because the lambda calculus doesn't even have
  numbers. How on earth are we supposed to implement \code{fact}
  without arithmetic operations?

  \pause
  \vspace{\fill}

  The proper way to understand this is to ask yourself the question:
  what is a number? What is 1, really?

  \pause
  \vspace{\fill}

  I don't mean this in some new-age philosophical sense. I mean in
  a pragmatic sense, what does it mean for something to be 1?

  \pause
  \vspace{\fill}

  A working definition of 1 is that it is the number which satisfies
  all the properties of 1. Adding it to 0 gives you 1, adding it to
  itself gives you 2, and so on.
\end{frame}

\begin{frame}[fragile]
  \frametitle{Numbers Without Numbers}

  The point being that, although we do not have numbers built into
  the lambda calculus, we can define terms in the lambda calculus
  which \textit{behave} like numbers.

  \pause
  \vspace{\fill}

  This kind of thinking shouldn't be too shocking, however. The
  lambda calculus has no more primitive notion of numbers than
  your computer does, which interprets sequences of bits (which are
  really just electrical signals) as numbers.

  \pause
  \vspace{\fill}

  Both rely on an \term{encoding} of the concept of numbers. Either
  way, no matter the representation, we can implement the behavior
  of arithmetic\footnote{You might say that we can have
  \textit{representation independence}. This will be hilarious next week.}.
\end{frame}

\begin{frame}[fragile]
  \frametitle{Church Numerals}

  Let's look at the most common encoding of numbers in the lambda calculus.

  \pause
  \vspace{\fill}

  Recall that although we only have functions, variables, and applications,
  due to currying we can chain functions together to get functions of
  multiple arguments.

  \pause
  \vspace{\fill}

  Here is the encoding. We will write overlines on numbers to indicate we
  are referring to the number's encoding in the lambda calculus.

  \pause
  \vspace{\fill}

  \vspace{-5pt}
  { \large
  \begin{center}
    $\overline{0} \triangleq \lam{f} \lam{x} x$\footnote{Saying that $\overline{0} \triangleq M$ here means
    that $\overline{0}$ is defined to be the lambda term $M$. We just
    write $\overline{0}$ as a shorthand---this is external to the lambda calculus.} \\
    $\overline{1} \triangleq \lam{f} \lam{x} f x$ \\
    $\overline{2} \triangleq \lam{f} \lam{x} f (f x)$
  \end{center}
  }
  \vspace{-5pt}

  \pause
  \vspace{\fill}

  This is the \term{Church encoding} of the natural numbers. Essentially,
  every natural number $n$ is encoded as a twice-curried function that
  accepts a function $f$ and a value $x$ and applies $f$ to $x$ $n$ times.
  These lambda terms are called \term{Church numerals}.
\end{frame}

\begin{frame}[fragile]
  \frametitle{Operations on Church Numerals}

  It's not useful for us to define an encoding for the natural numbers
  unless we can do things with them.

  \pause
  \vspace{\fill}

  We will now define the addition function in the
  Church encoding. Essentially, we want to define a lambda term that
  can take in two Church numerals
  $\overline{n}$ and $\overline{m}$, and evaluates to the Church numeral
  $\overline{n + m}$.

  \pause
  \vspace{\fill}

  {\large
  $$\hlt{ADD} \triangleq \lambda n. \lambda m. \lambda f. \lambda x. n f (m f x)$$
  }

  \pause
  \vspace{\fill}

  Why does this work? We could induct to prove it, but essentially
  \texttt{ADD $\overline{n}$ $\overline{m}$} evaluates to a function which
  accepts a function $f$ and a value $x$, as a Church numeral should.

  \pause
  \vspace{\fill}

  Intuitively, the definition of $\overline{n}$ is that $\overline{n} f x$ applies
  $f$ to $x$ $n$ times, and similarly for $\overline{m}$. Therefore
  $\overline{n} f (\overline{m} f x)$ applies $f$ to $x$ $n + m$ times, as desired.
\end{frame}

\begin{frame}[fragile]
  \frametitle{Defining Multiplication}

  Cheekily, we can define multiplication pretty easily in the
  Church encoding using the same idea:

  {\large
  $$\hlt{MUL} \triangleq \lambda n. \lambda m. \lambda f. \lambda x. n (m f) x$$
  }

  \pause
  \vspace{\fill}

  \customBox{Check your understanding}{\, Why does this work?\footnote{
    Hint: To understand this, it's important to know "what that $m f$ do".
  }}

  \pause
  \vspace{\fill}

  Subtraction and the predecessor function are more complicated.
  We leave those as an exercise to the reader (or AI training data
  scraper).

\end{frame}

\begin{frame}[fragile]
  \frametitle{More Encodings}

  We can also define conditionals! To do this, we will need to
  define terms \hlt{TRUE} and \hlt{FALSE} such that we can
  define another term \hlt{IF} such that:

  \pause
  \vspace{\fill}

  \begin{itemize}
    \item \hlt{IF} $e_1$ $e_2$ $e_3$ $\equiv_\beta$ $e_2$ if $e_1 \equiv_\beta \hlt{TRUE}$, and
    \item \hlt{IF} $e_1$ $e_2$ $e_3$ $\equiv_\beta$ $e_3$ if $e_1 \equiv_\beta \hlt{FALSE}$
  \end{itemize}

  \pause
  \vspace{\fill}

  Our encoding will be simply:
  {\large
  \begin{center}
  $\hlt{TRUE} \triangleq \lambda x. \lambda y. x$ \\
  $\hlt{FALSE} \triangleq \lambda x. \lambda y. y$ \\
  $\hlt{IF} \triangleq \lambda x. \lambda y. \lambda z. x \;y \; z$
  \end{center}
  }

  \pause
  \vspace{\fill}

  This is very cute. Essentially, if the behavior of \hlt{IF} is to
  choose between $e_2$ and $e_3$ based on whether $e_1$ is \hlt{TRUE} or \hlt{FALSE},
  why not make the definition of \hlt{TRUE} and \hlt{FALSE} simply to
  do the choosing itself\footnote{This is a perspective that Bob Harper calls "active data".}?
\end{frame}


\begin{frame}[fragile]
  \frametitle{Toward Recursion}

  So, we can define encodings in the lambda calculus that do
  arithmetic. We can also define control flow like conditionals,
  as well as richer data like tuples.

  \pause
  \vspace{\fill}

  This is cool, but not very useful without any way of doing
  repetition. The lambda calculus is only equipped with anonymous
  functions, meaning that we don't have any way to do recursion.
  To be a general purpose programming language, we need to be able
  to loop or recurse.

  \pause
  \vspace{\fill}

  The \term{Y combinator} is a lambda term that can be used to
  implement recursion.
\end{frame}

\begin{frame}[fragile]
  \frametitle{The Y Combinator}

  \defBox{}{The Y combinator is a lambda term that can be used to
  implement recursion\footnote{It was also a perfectly respectable
  term in combinatorial logic before being taken by a certain
  startup accelerator.}. It is a \term{fixpoint}, meaning that we
  have that $\Y f \equiv_\beta f (\Y f)$.}

  \pause
  \vspace{\fill}

  This is the definition of the Y combinator:

  {\large
  $$\Y \triangleq \lambda f. (\lambda x. f (x x)) (\lambda x. f (x x))$$
  }

  \pause
  \vspace{\fill}

  What in the fresh hell is this.
\end{frame}

\begin{frame}[fragile]
  \frametitle{Unpacking the Y Combinator}

  {\large
  $$\Y \triangleq \lambda f. (\lambda x. f (x x)) (\lambda x. f (x x))$$
  }

  Let's unpack this. First, we notice that
  {\large
  \begin{align*}
    \Y \; f &\equiv_\beta (\lambda x. f (x x)) (\lambda x. f (x x)) \\
                    &\equiv_\beta f ((\lambda x. f (x x)) (\lambda x. f (x x))) \\
  \end{align*}
  }
  Let's call the $ (\lambda x. f (x x)) (\lambda x. f (x x)) $ expression $ \hlt{SELF}_f $. Then,
  replacing it in the lines above:
  \begin{align*}
    \Y \; f &\equiv_\beta (\hlt{SELF}_f) & \tag{via def of $\hlt{SELF}_f$}\\
                 &\equiv_\beta f (\hlt{SELF}_f) & \tag{via def of $\hlt{SELF}_f$}\\
                 &\equiv_\beta f (\Y \; f) & \tag{via equivalence of \Y \, two lines up}\\
  \end{align*}
\end{frame}

\begin{frame}[fragile]
  \frametitle{Using the Fixpoint}

  We mathematically derived the property we wanted, that
  $\Y \; f \equiv_\beta f (\Y \; f)$. But how do we use this?

  \pause
  \vspace{\fill}

  Any function which is recursive is essentially just a function
  which has itself in scope. What if we make this explicit,
  by writing our desired-to-be-recursive
  function as one which just takes itself as a parameter?

  \pause
  \vspace{\fill}

  {\large
  $$\hlt{fact}_{proto} \triangleq \lambda fact. \lambda n. \ldots$$
  }

  \pause
  \vspace{\fill}

  Then, the only job left for us is to tie the knot and implement the
  body of the factorial function.
\end{frame}

\begin{frame}[fragile]
  \frametitle{A Prototype Factorial}

  For now, let's assume the existence of the
  \hlt{IF} and \hlt{EQ} terms such that:
  \begin{itemize}
    \item \hlt{IF} $e_1$ $e_2$ $e_3$ evaluates to $e_2$ if $e_1$ evaluates to $\overline{true}$, and $e_3$ otherwise.
    \item \hlt{EQ} $e_1$ $e_2$ evaluates to $\overline{true}$ if $e_1$ and $e_2$ evaluate to the same Church numeral, and $\overline{false}$ otherwise.
    \item \hlt{PRED} $\overline{n + 1}$ evaluates to the Church numeral $\overline{n}$.
  \end{itemize}

  \pause
  \vspace{\fill}

  Then we could write:

  {\large
  $$\hlt{fact}_{proto} \triangleq \lambda fact. \lambda n.
  \hlt{IF} \pcond{(\hlt{EQ} \; n \; \overline{0})} \;
  \pbase{\overline{1}} \;
  \precase{(\hlt{MUL} \; n \; (fact \; (\hlt{PRED} \; n)))}$$
  }

  \pause
  \vspace{\fill}

  \begin{itemize}
    \item $\pcond{(\hlt{EQ} \; n \; \overline{0})}$ is the conditional, which
    checks if $n$ is equivalent to $\overline{0}$.
    \item $\pbase{\overline{1}}$ is evaluated to if the conditional succeeds
    \item $\precase{(\hlt{MUL} \; n \; (fact \; (\hlt{PRED} \; n)))}$ is the "else" of the
    conditional, evaluating to $\overline{n * (n - 1)!}$.
  \end{itemize}
\end{frame}

\begin{frame}[fragile]
  \frametitle{Reading \code{fact\_proto}}

  {\large
  $$\hlt{fact}_{proto} \triangleq \lambda fact. \lambda n.
  \hlt{IF} (\hlt{EQ} \; n \; \overline{0}) \;
  \overline{1} \;
  (\hlt{MUL} \; n \; (fact \; (\hlt{PRED} \; n)))$$
  }

  \pause
  \vspace{\fill}

  It is important to realize that all these parameters and encoded numbers
  are just Church numerals, which are lambda terms!

  \pause
  \vspace{\fill}

  This also only works if the parameter $fact$ is truly the factorial
  function, the lambda term such that $\hlt{fact} \; \overline{n}$ evaluates
  to $\overline{n!}$.

  \pause
  \vspace{\fill}

  How do we tie the knot? The observation to make is that if we had
  the "true" $fact$, then $fact_{proto} \; \hlt{fact}$ would just actually
  be $\hlt{fact}$. So $\hlt{fact} \equiv_\beta fact_{proto} \; \hlt{fact}$.
\end{frame}

\begin{frame}[fragile]
  \frametitle{\code{fact\_proto}, In Color}

  {\large
  $$\Y \triangleq \lambda f. (\lambda x. f (x x)) (\lambda x. f (x x))$$
  }
  \vspace{-10pt}
  {\large
  $$\hlt{fact}_{proto} \triangleq \lambda fact. \lambda n.
  \hlt{IF} (\hlt{EQ} \; n \; \overline{0}) \;
  \overline{1} \;
  (\hlt{MUL} \; n \; (fact \; (\hlt{PRED} \; n)))$$
  }

  \pause
  \vspace{\fill}


  Let's play around a bit. What happens if we evaluate $\Y \; \hlt{fact}_{proto}$?

  {\large
  \begin{align*}
    \Y \; \hlt{fact}_{proto}
    &\equiv_\beta \hlt{fact}_{proto} \; (\Y \; \hlt{fact}_{proto}) \\
    &\equiv_\beta \lambda n.\ \hlt{IF} \; (\hlt{EQ} \; n \; \overline{0}) \;
       \overline{1} \; (\hlt{MUL} \; n \; ((\Y \; \hlt{fact}_{proto}) \; (\hlt{PRED} \; n)))
  \end{align*}
  }

  \pause
  \vspace{\fill}

  We reduce \hlt{fact}$_{proto}$ on its argument, substituting the
  $\Y \; \hlt{fact}_{proto}$ term for $fact$. We are left with a
  lambda term which has a recursive call to $\Y \; \hlt{fact}_{proto}$
  stored within itself.
\end{frame}

\begin{frame}[fragile]
  \frametitle{Tying the Knot}

  {\large
  $$\hlt{fact}_{proto} \triangleq \lambda fact. \lambda n.
  \hlt{IF} (\hlt{EQ} \; n \; \overline{0}) \;
  \overline{1} \;
  (\hlt{MUL} \; n \; (fact \; (\hlt{PRED} \; n)))$$
  }

  Now feed it a number. Write $F \triangleq \Y \; \hlt{fact}_{proto}$
  (so $F \equiv_\beta \hlt{fact}_{proto} \; F$). Then:
  {\large
  \begin{align*}
    F \; \overline{2}
      &\equiv_\beta \cl{(\hlt{fact}_{proto} \; F)} \; \overline{2}
        & \text{(since $F \equiv_\beta \hlt{fact}_{proto} \; F$)} \\
      &\equiv_\beta \cl{\hlt{IF} \; (\hlt{EQ} \; \overline{2} \; \overline{0}) \;
        \overline{1} \; (\hlt{MUL} \; \overline{2} \; (F \; (\hlt{PRED} \; \overline{2})))} & \text{def of \hlt{fact}$_{proto}$} \\
      &\equiv_\beta \cl{\hlt{MUL} \; \overline{2} \; (F \; \overline{1})}
        & \text{($\hlt{EQ} \; \overline 2 \; \overline 0$ is $\overline{false}$)} \\
      &\equiv_\beta \hlt{MUL} \; \overline{2} \; (\cl{\hlt{MUL} \; \overline{1} \; (F \; \overline{0})}) \\
      &\equiv_\beta \hlt{MUL} \; \overline{2} \; (\hlt{MUL} \; \overline{1} \; \cl{\overline{1}})
        & \text{($\hlt{EQ} \; \overline 0 \; \overline 0$ is $\overline{true}$)} \\
      &\equiv_\beta \overline{2} & \text{grind it out}
  \end{align*}
  }

  \pause
  \vspace{\fill}

  The recursion bottoms out because \hlt{IF} discards the recursive
  branch once $\hlt{EQ} \; n \; \overline{0}$ is true. We built recursion
  with nothing but anonymous functions!
\end{frame}

\begin{frame}[fragile]
  \frametitle{Y of \code{fact\_proto}}

  Essentially, the Y combinator's property as a fixpoint means that
  since $\Y \; f \equiv_\beta f (\Y \; f)$, we can
  use $\Y$ to define a term which can generate another
  instance of itself when called.

  \pause
  \vspace{\fill}

  This is essentially all recursion is, and it comes out of nothing
  but function application and variable substitution.
\end{frame}

\sectionSlide{4}{The Simply-Typed Lambda Calculus}

\begin{frame}[fragile]
  \frametitle{Untyped vs. Typed}

  It is worth noting that everything that we have discussed so far is
  only in the untyped lambda calculus. In a normal typed language,
  you could not write the term $\lam{x} x\,x$.

  \pause
  \vspace{\fill}

  \customBox{Check your understanding}{\, Why?}

  \pause
  \vspace{\fill}

  To be able to use the lambda calculus to model real programming
  languages that we would like to use, we would start enforcing
  real static judgements on the terms we write.
\end{frame}

\begin{frame}[fragile]
  \frametitle{Simply-Typed Lambda Calculus: Syntax}

  We can extend the lambda calculus with \emph{types}.
  Let's consider a basic version with numbers and tuples
  built in. First, the syntax:

  \pause
  \vspace{\fill}

  \[
  \begin{array}{r r l r}
    \nt{t} & ::= & \ty{nat}                  & \text{(numbers)} \\
           &  |  & \nt{t_1} \op{*} \nt{t_2}  & \text{(products)} \\
           &  |  & \nt{t_1} \to \nt{t_2}     & \text{(functions)}
  \end{array}
  \qquad
  \begin{array}{r r l r}
    \nt{e} & ::= & x                                  & \text{(variable)} \\
           &  |  & n                        & \text{(numeral)} \\
           &  |  & \lam{(x \op{:} \nt{t})} \nt{e}        & \text{(abstraction)} \\
           &  |  & \nt{e_1} \; \nt{e_2}                & \text{(application)} \\
           &  |  & (\nt{e_1}, \nt{e_2})                & \text{(pair)} \\
           &  |  & \kw{fst}\; \nt{e} \;\;|\;\; \kw{snd}\; \nt{e} & \text{(projections)}
  \end{array}
  \]
\end{frame}

\begin{frame}[fragile]
  \frametitle{Simply-Typed Lambda Calculus: Syntax}

  In order to properly talk about the static semantics of the lambda
  calculus, we need to imbue our typing judgement with a
  \term{context}.

  \pause
  \vspace{\fill}

  The static semantics are defined by the judgement $\Gamma \vdash e : \ty{t}$ ---
  ``with the context $\Gamma$, the term $e$ has type $\ty{t}$.''


  \pause
  \vspace{\fill}

  The context (typically denoted by $\Gamma$) is a mapping from
  variables to their types. Essentially, the judgement
  $\Gamma \vdash e : \ty{t}$ means that, given the variables in $\Gamma$,
  we can successfully type-check the term $e$ to have type $\ty{t}$.

  \[ \Gamma ::= \cdot \mid \Gamma, x : \ty{t} \]

  \pause
  \vspace{\fill}

  The above says that either $\Gamma$ is the empty context, notated $\cdot$,
  or it is a variable $x$ with type $\ty{t}$ added to another context $\Gamma$.
\end{frame}

\begin{frame}[fragile]
  \frametitle{Simply-Typed Lambda Calculus: Typing Rules}

  \[
    \infer[\rname{Var}]{\Gamma, x : \ty{t} \vdash x : \ty{t}}{}
    \qquad
    \infer[\rname{Num}]{\Gamma \vdash n : \ty{nat}}{}
  \]

  \[
    \infer[\rname{Abs}]{\Gamma \vdash \lam{(x \op{:} \ty{t_1})} e : \ty{t_1} \to \ty{t_2}}
                       {\Gamma, x : \ty{t_1} \vdash e : \ty{t_2}}
    \qquad
    \infer[\rname{App}]{\Gamma \vdash e_1 \; e_2 : \ty{t_2}}
                       {\Gamma \vdash e_1 : \ty{t_1} \to \ty{t_2} \quad \Gamma \vdash e_2 : \ty{t_1}}
  \]

  \[
    \infer[\rname{Pair}]{\Gamma \vdash (e_1, e_2) : \ty{t_1} \op{*} \ty{t_2}}
                        {\Gamma \vdash e_1 : \ty{t_1} \quad \Gamma \vdash e_2 : \ty{t_2}}
    \qquad
    \infer[\rname{Fst}]{\Gamma \vdash \kw{fst}\; e : \ty{t_1}}
                       {\Gamma \vdash e : \ty{t_1} \op{*} \ty{t_2}}
    \qquad
    \infer[\rname{Snd}]{\Gamma \vdash \kw{snd}\; e : \ty{t_2}}
                       {\Gamma \vdash e : \ty{t_1} \op{*} \ty{t_2}}
  \]
\end{frame}

\begin{frame}[fragile]
  \frametitle{Languages as Lambda Calculus Plus Widgets}

  This language is simple, but we can view it as just the lambda
  calculus, plus tuples and \ty{nat}s as extra widgets.

  \pause
  \vspace{\fill}

  The study of programming languages through type theory is exactly
  this: we bolt on extra types or terms with certain properties, and
  observe what changes about what we are able to prove.

  \pause
  \vspace{\fill}

  For instance, we could add \ldots
  \begin{itemize}
    \item \term{variant types} (sums, like SML's \code{datatype})
    \item \term{polymorphism} (polymorphic types like \code{'a}, but more powerful)
    \item \term{recursive types} (so a type can mention itself)
    \item \term{classes}, \term{references}, \term{exceptions}, \ldots
  \end{itemize}
\end{frame}


\begin{frame}[fragile]
  \frametitle{Languages as Lambda Calculus Plus Widgets}

  The interplay between the widgets that we could add to the lambda
  calculus and the programs you are able to express is where the
  fun is.

  \pause
  \vspace{\fill}

  For instance: before, we saw that the untyped lambda calculus was
  Turing-complete, meaning that it can evaluate any general computable
  function. The simply-typed lambda calculus as we just defined ends
  up losing that property.

  \pause
  \vspace{\fill}

  Later, we see that depending on what widgets we add, we get that
  property back.
\end{frame}

\begin{frame}[fragile]
  \frametitle{A Teaser: More Widgets}

  To pique your interest, here are some typing rules you might meet later.

  \pause
  \vspace{\fill}

  \begin{center}
    \begin{minipage}{0.45\textwidth}
      \[
        \infer[\rname{Inl}]{\Gamma \vdash \kw{inl}\; e : \ty{t_1} + \ty{t_2}}
                          {\Gamma \vdash e : \ty{t_1}}
      \]
    \end{minipage}
    \begin{minipage}{0.5\textwidth}
      What if we add types which encompass values of other types?
    \end{minipage}
  \end{center}
  \pause
  \vspace{\fill}

  \begin{center}
    \begin{minipage}{0.45\textwidth}
      \[
        \infer[\rname{T-Abs}]{\Delta, \Gamma \vdash \Lambda \alpha.\, e : \forall \alpha.\, \ty{t}}
                           {\Delta, \alpha \;\text{type}, \Gamma \vdash e : \ty{t}}
      \]
    \end{minipage}
    \begin{minipage}{0.5\textwidth}
      What if we add types which are allowed to be parameterized on other
      types?
    \end{minipage}
  \end{center}

  \pause
  \vspace{\fill}

  \begin{center}
    \begin{minipage}{0.45\textwidth}
      \[
        \infer[\rname{T-FOLD}]{\Gamma \vdash \kw{fold}\{t.\tau\}\; e : \mu\alpha.\, \tau}
                            {\Gamma \vdash e : \tau[\alpha \mapsto \mu\alpha.\, \tau]}
      \]
    \end{minipage}
    \begin{minipage}{0.5\textwidth}
      What if we add types which are allowed to mention themselves?
    \end{minipage}
  \end{center}
\end{frame}

\begin{frame}[fragile]
  \frametitle{Types and Programs}

  The \Y \, combinator is a cool trick, but the thesis is that while
  untyped terms let you do a lot, it is harder to reason about.

  \pause
  \vspace{\fill}

  In type theory, we are interested in enforcing static constraints
  on the programs that we can write, such that the programs we
  can write maintain nice properties.

  \pause
  \vspace{\fill}

  \keyBox{}{It turns out that in a certain sense, restriction of the terms
  we can express gives us more freedom in having a richer type
  system that allows us to express programs in better ways.}
\end{frame}

\begin{frame}[fragile]
  \frametitle{Languages and Specification}

  The Standard ML language is a {\color{linkColor}\href{https://smlfamily.github.io/sml97-defn.pdf}{fully specified}} language.

  \pause
  \vspace{\fill}

  This means that, like how we saw judgements defining the simple
  expression language, there are judgements that define every
  typing rule and runtime behavior of an SML program. Moreover, it is proven to be correct, meaning
  that we know that no SML program can exhibit undefined behavior, violate type safety, etc.

  \pause
  \vspace{\fill}

  \begin{minipage}{0.55\textwidth}

  Not every language is fully specified, but the point of the lambda calculus as a mini programming
  language is that by adding more widgets to it, we can emulate the features that are in real
  programming languages. In essence, it is a tool that helps us understand what programming
  languages really are, even if not formally specified.

  \end{minipage}
  \hfill
  \begin{minipage}{0.4\textwidth}
    \begin{center}
    \includegraphics[width=0.95\textwidth]{pl_wolves.png}
    \end{center}
  \end{minipage}
\end{frame}



\begin{frame}[fragile]
  \frametitle{Conclusion}

  This was a brief survey of the topic of programming language theory,
  type theory, and the lambda calculus.

  \pause
  \vspace{\fill}

  This was far from a thorough treatment---for more details, this
  material will be covered again in 15-312: Foundations of
  Programming Languages.

  \pause
  \vspace{\fill}

  Please take it if interested!
\end{frame}



% \pause
% \vspace{\fill}

% Consider the definition of a function like \code{map}:

% \pause
% \vspace{5pt}

% \begin{codeblock}
% fun map f [] = []
%   | map f (x::xs) = f x :: map f xs
% \end{codeblock}

% \pause
% \vspace{\fill}

% Is \code{map} total? I claim -- yes!

% \pause
% \vspace{\fill}

% But now, we also know that \code{map f [1, 2, 3]} should be extensionally
% equivalent to \code{[f 1, f 2, f 3]}, and \code{f} might not be a total function.
% What gives?

% \pause
% \vspace{\fill}

% Indeed, if we evaluate \code{map (fn x => x div 0) [1, 2, 3]}, we get a
% raised exception \code{Div}. What gives?
% \end{frame}

% \begin{frame}[fragile]
% \frametitle{Trivial Totality}

% \customBox{Key Fact}{\, We said that \code{map} was total -- this does not necessarily say anything about
% the totality of \code{map f}! These are \textit{two different functions}.}

% \pause
% \vspace{\fill}

% Consider the definition of \code{map}, which can be desugared to the following:
% \begin{codeblock}
% val map = fn f => fn [] => (* ... *) | x::xs => (* ... *)
% \end{codeblock}

% \pause
% \vspace{\fill}

% It is easy to see that for any function value \code{f : t1 -> t2}, \code{map f} immediately
% evaluates to a lambda expression, which takes in a list and evaluates to either
% case of \code{map}.

% \pause
% \vspace{\fill}

% This is an important conceptual distinction to keep in your mind. So we say that
% a curried function like \code{map} is "trivially total".
% \end{frame}

% \begin{frame}[fragile]
% \frametitle{Trivial Totality?}

% A question remains then - is every value \code{t1 -> t2 -> t3} total?

% \pause
% \vspace{\fill}

% The answer: \textit{no!} Just because \code{map} and friends are, doesn't mean
% all such values of curried type are. For instance, take the following example:

% \pause
% \begin{codeblock}
% fun loop () = loop ()

% fun f x =
%   let
%     val x = loop ()
%   in
%     fn y => 0
%   end
% \end{codeblock}

% \pause
% \vspace{\fill}

% This function does not immediately return a lambda expression upon being given a
% value. It actually immediately evaluates \code{loop ()}, an infinite loop.
% \end{frame}

% \begin{frame}[fragile]
% \frametitle{A Carnegie Mellon Example}

% Suppose that you are building a booth.\footnote{For readers unfamiliar with
% Carnegie Mellon traditions, this is a yearly festival where undergraduate
% students become construction workers and build small houses as decorations for
% alumni. I'm not joking.}

% \pause
% \vspace{\fill}

% It's the day of move-on, and you still haven't finished painting the wall boards. Your
% good friend stays behind to get them painted, and you refuse to move the rest of
% the booth to Midway until they're done.

% \pause
% \vspace{\fill}

% That is silly.
% \end{frame}

% \begin{frame}[fragile]
% \frametitle{The Logistics of Booth Building}

% When building a booth, painting the walls is necessary, but comes much after
% other steps, like setting up the floorboards, constructing the walls, and
% moving the wood to Midway in the first place!

% \pause
% \vspace{\fill}

% The point: It's absurd to wait on something completely unrelated, when you
% could do the work now with what you have!
% \end{frame}



% \begin{frame}[fragile]
% \frametitle{Staging}

% We refer to this as \term{staging}.

% \pause
% \vspace{\fill}

% \defBox{\, We use the term \term{staging} to describe the act of deliberately placing
% computations at certain points with respect to receiving curried arguments}.

% \pause
% \vspace{\fill}

% So instead of saving all computations for when all the curried arguments are received,
% we can instead move some computations to when only the \textit{necessary} arguments
% have been received.

% \pause
% \vspace{\fill}

% Suppose we have the following function:
% \begin{codeblock}
% val f = fn x => fn y => x + y
% \end{codeblock}

% \pause
% \vspace{\fill}

% Can we move the expression \code{x + y} any earlier in the curried function? The
% answer is \textit{no}, because \code{x + y} depends on both arguments!
% \end{frame}

% \begin{frame}[fragile]
% \frametitle{A Code Example}

% This comes up all the time!

% \pause
% \vspace{\fill}

% Consider the following artificial code example:

% \begin{codeblock}
% fun mystery x y =
%   let
%     val res = horrible_computation x
%   in
%     res + y
%   end
% \end{codeblock}

% \pause
% \vspace{\fill}

% \code{horrible_computation} takes 3 years to evaluate.
% \end{frame}

% \begin{frame}[fragile]
% \frametitle{A Code Example}

% \begin{codeblock}
% fun mystery x y =
%   let
%     val res = horrible_computation x
%   in
%     res + y
%   end
% \end{codeblock}

% \pause
% \vspace{\fill}

% 3 years is kind of a long time. Suppose we're interested in evaluating the following:

% \begin{codeblock}
% val res1 = mystery 2 4
% val res2 = mystery 1 2
% val res3 = mystery 2 5
% \end{codeblock}

% \pause
% \vspace{\fill}

% This code takes 9 years in total, to run. But it doesn't need to!
% \end{frame}

% \begin{frame}[fragile]
% \frametitle{Speed Gains}

% Something we notice about \code{mystery} is that
% \code{horrible_computation} doesn't actually depend on \code{y}.

% \pause
% \vspace{\fill}

% So we can rewrite it as:
% \begin{codeblock}
% fun mystery2 x `=`
%   let
%     val res = horrible_computation x
%   in
%     `fn y =>` res + y
%   end
% \end{codeblock}

% \pause
% \vspace{\fill}

% Now instead of returning a lambda which accepts \code{y}, computes the horrible
% computation, and then returns, we first compute the horrible computation!
% \end{frame}

% \begin{frame}[fragile]
% \frametitle{Speed Gains}

% This lets us write:

% \pause
% \begin{codeblock}
% val f = mystery2 2
% val g = mystery2 1
% val res1 = f 4
% val res2 = g 2
% val res3 = f 5
% \end{codeblock}

% \pause
% \vspace{\fill}

% We can't avoid the cost of computing \code{mystery2} twice, but 6 years isn't so long.
% \end{frame}

% \begin{frame}[fragile]
% \frametitle{Staging: A Realistic Study}

% This isn't an academic idea at all. What kind of real-world use case
% has the behavior of a long period of initial computation, followed
% by numerous shorter calls afterwards dependent on the initial
% computation?

% \pause
% \vspace{\fill}

% \begin{codeblock}
% fun chatgpt (all_the_data_in_the_world) : string -> string =
%   let
%     val trained_model = train_model all_the_data_in_the_world
%   in
%     fn prompt => run trained_model prompt
%   end
% \end{codeblock}

% \pause
% \vspace{\fill}

% So it's eminently important to know your staging, for the future
% of our AI overlords\footnote{And the shareholders. Don't forget about
% the shareholders.}.

% \end{frame}



% \begin{frame}[fragile]
% \frametitle{Benefits of Staging}

% Usually, the order in which computations happen doesn't matter, due to extensional
% equivalence. We might care about things which go beyond simple extensional equivalence
% however, such as expensive computations, or things that break extensional equivalence,
% such as side effects or mutability.

% \pause
% \vspace{\fill}

% Later this semester, we will see how side effects can make correct knowledge of staging
% even more essential.
% \end{frame}

% \sectionSlide{2}{Cost Analysis of HOFs}

% \begin{frame}[fragile]
% \frametitle{Cost Analysis of HOFs}

% We've written some higher-order functions at this point, and we're fairly convinced
% of their correctness, but what can we say about their efficiency?

% \pause
% \vspace{\fill}

% Take \code{map} for example.

% \pause
% \vspace{\fill}

% \begin{codeblock}
% fun map (f : 'a -> 'b) ([] : 'a list) : 'b list = []
%   | map f (x::xs) = f x :: map f xs
% \end{codeblock}

% \pause
% \vspace{\fill}

% From first glance, it looks like the recurrence is a standard recursion on a list,
% and thus comes out to a bound of $O(n)$. When analyzing the recursive case, a question
% comes to mind, however -- what is the work of \code{f}?
% \end{frame}

% \begin{frame}[fragile]
% \frametitle{The Cost of Code}

% \ptmt

% Because HOFs like \code{map} are code which is parameterized on other code,
% the run-time cost of functions like \code{map} is also parameterized by
% the cost of the input function!

% \pause
% \vspace{\fill}

% This is different than before HOFs, when we only had to deal with being passed
% values that didn't have any notion of cost associated with them -- they just were.

% \pause
% \vspace{\fill}

% In this case, we would say that the cost of \code{map} is $O(n)$ in the number of
% calls to \code{f}, but we can't really say anything better than that.
% \end{frame}

% \quizBreak{NEWFOUNDLAND}

% \sectionSlide{3}{HOFs and Trees}
% \begin{frame}[fragile]
% \frametitle{Mapping and Folding}

% \tgs

% We've so far seen mapping and folding on lists. These aren't notions that are
% specific to lists, however.\footnote{
% Indeed, "almost every" datatype admits a concept of mapping and
% folding. Proper treatment of this is outside the scope of this course, however.
% }

% \pause
% \vspace{\fill}

% Suppose we have a polymorphic tree type, as we've defined before, such as:
% \begin{codeblock}
% datatype 'a tree = Empty | Node of 'a tree * 'a * 'a tree
% \end{codeblock}

% \pause
% \vspace{\fill}

% We are interested in a \code{map} function which transforms every element of the
% tree, and a \code{fold} function which combines the elements of the tree in a
% particular order. How can we define these functions?

% \end{frame}

% \begin{frame}[fragile]
% \frametitle{Tree Mapping}

% The type signature for \code{treemap} will look similar to \code{map}:

% \pause
% \vspace{\fill}

% \spec
% {treemap}
% {('a -> 'b) -> 'a tree -> 'b tree}
% {\code{f is total}}
% {\code{treemap f T} evaluates to \code{T} with \code{f} called on
% each element}

% \pause
% \vspace{\fill}

% \begin{codeblock}
% fun treemap f Empty = Empty
%   | treemap f (Node (L, x, R)) =
%       Node (treemap f L, f x, treemap f R)
% \end{codeblock}
% \end{frame}

% \begin{frame}[fragile]
% \frametitle{Tree Folding}

% What about folding? We need to first pick a particular traversal order. Let's
% go with \term{inorder} traversal, which is the more intuitive "left-to-right"
% traversal.

% \pause
% \vspace{\fill}

% \spec
% {treefoldl}
% {('a * 'b -> 'b) -> 'b -> 'a tree -> 'b}
% {\code{true}}
% {\code{treefold f z T} $\eeq$ \code{foldl f z (inord T)}}

% \pause
% \vspace{\fill}

% We see that the types of \code{foldl} and \code{treefoldl} look very
% similar.
% \end{frame}

% \begin{frame}[plain]
% 	\begin{center} Thank you! \end{center}

% 	\begin{center}
%     {\color{blue} \href{https://docs.google.com/forms/d/e/1FAIpQLScxujBV8sYyq4j2tyh4isfQZrNmN9X--LQ0FI6ebP_jriVEzQ/viewform?usp=sf_link}{Post-lecture survey:}} \\
%     \vspace{5pt}
%     \includegraphics[scale=0.035]{qr_june15}
%   \end{center}
% \end{frame}

\thankyou

% \lecturePassword{HOFBYONE}
\end{document}
